\documentclass[12pt,a4paper]{article}
% Preamble bersama untuk kedua laporan TK 2 (soal1 dan soal2).
% Mengikuti "Petunjuk Penulisan Dokumen Project Anum":
%   - A4, font standar LaTeX (Latin Modern), 12pt, margin >= 2.54 cm
%   - judul, rangkuman, daftar isi TIDAK dinomori; bagian lain dinomori dari 1
%   - tabel/gambar/kode/persamaan dinomori mengikuti bagian: Tabel 3.1, Gambar 3.2, Kode 3.1
\usepackage[T1]{fontenc}
\usepackage[utf8]{inputenc}
\usepackage[indonesian]{babel}
\usepackage{lmodern}
\usepackage[a4paper,margin=2.54cm]{geometry}
\usepackage{amsmath,amssymb,booktabs,array,graphicx,float,microtype}
\usepackage{algorithm,algpseudocode}
\usepackage{listings}
\usepackage[hidelinks]{hyperref}
\emergencystretch=1em

\numberwithin{equation}{section}
\counterwithin{table}{section}
\counterwithin{figure}{section}
\counterwithin{algorithm}{section}
\floatname{algorithm}{Kode}
\renewcommand{\lstlistingname}{Kode}
\algrenewcommand\algorithmicrequire{\textbf{Masukan:}}
\algrenewcommand\algorithmicensure{\textbf{Keluaran:}}
\setlength{\belowcaptionskip}{6pt}
\lstset{language=Python, basicstyle=\ttfamily\footnotesize, breaklines=true,
  columns=fullflexible, keepspaces=true, showstringspaces=false, frame=single}

\newcommand{\R}{\mathbb{R}}
\newcommand{\T}{\mathsf{T}}
\newcommand{\norm}[1]{\left\lVert #1\right\rVert_2}
\newcommand{\bx}{\mathbf{x}}
\newcommand{\bd}{\mathbf{d}}
\newcommand{\bg}{\mathbf{g}}
\newcommand{\bs}{\mathbf{s}}
\newcommand{\by}{\mathbf{y}}

% \input biasa di dalam tabular gagal bila diikuti \bottomrule; pakai primitif \@@input
\makeatletter\newcommand{\isitabel}[1]{\@@input #1 }\makeatother

% Grafik dicari di report/figures/ dan results/ (relatif terhadap folder report/)
\graphicspath{{figures/}{../results/}}

\begin{document}

% ---- Judul (1 baris, tidak dinomori) ----
\begin{center}
{\Large\bfseries Soal 1: Optimisasi Fungsi Rosenbrock dengan CGM dan Quasi-Newton}\par\medskip
Kelompok A9 -- TK 2 Analisis Numerik
\end{center}

% ---- Rangkuman (~1/2 halaman, tidak dinomori). Ditulis TERAKHIR. ----
\section*{Rangkuman}
% TODO: masalah, metode (CGM, quasi-Newton, pencarian langkah), hasil utama (iterasi, evaluasi, akurasi), kesimpulan.

\setcounter{tocdepth}{1}
\tableofcontents
\clearpage

% ---- Bagian bernomor ----
\section{Pendahuluan}
% TODO (ditulis terakhir): deskripsi masalah, tujuan, cara studi dilakukan.

\section{Metode Pencarian Langkah}\label{sec:linesearch}
% Tugas (i): binary search vs Armijo; Tabel hasil uji; pilihan metode + alasan.
% Pseudocode (Kode 2.1, 2.2) -- beri caption + penjelasan di teks.

\section{Varian Algoritma}\label{sec:varian}
% Tugas (ii): dua rumus beta_k untuk CGM, dua rumus pembaruan untuk quasi-Newton.

\section{Perbandingan Hasil}\label{sec:hasil}
% Tugas (iii): tabel 8 baris (uji x algoritma x varian), jarak ke x* = (1,1).

\section{Visualisasi}\label{sec:visual}
% Tugas (iv): lintasan pada kontur f (skala log) + grafik log ||grad f|| vs k. Jelaskan cara menilai kesetaraan akurasi.

\section{Analisis}\label{sec:analisis}
% Tugas (v): jawab 4 pertanyaan secara singkat.

\section{Kesimpulan}
% TODO: 1--2 paragraf.

\begin{thebibliography}{9}
\renewcommand{\refname}{Referensi}
\bibitem{nocedal} J. Nocedal dan S. J. Wright, \emph{Numerical Optimization}, edisi ke-2. Springer, 2006.
% TODO: tambah referensi yang benar-benar dipakai (bukan slide kuliah).
\end{thebibliography}

\appendix
\section{Lampiran}
% Kode program (bukan pseudocode), tabel rinci. Contoh:
% \lstinputlisting[caption={Implementasi CGM.},label={code:cgm}]{../code/cgm.py}
\end{document}
